\section{The full support lattice and the arithmetic site's coefficient extension}
\label{sec:arithmetic-site-lattice-map}

The human source used here is Alain Connes and Caterina Consani,
\href{https://arxiv.org/abs/1502.05580}{\emph{Geometry of the Arithmetic
Site}}, original author source \texttt{arithmeticsite\_Adv\_final1.tex},
2074 lines, SHA--256
\par\noindent
\texttt{bd9da2b9e3c8846e306af18dbf745d7cf603450d6630a00e239da4235a7320b4}.
\par\noindent
The source mathematics actually read for this calculation is lines
331--407, 440--503, 598--704, and 850--1010. These ranges include the
definition of the characteristic-one semirings and Frobenius,
the rank-one description of points, the arithmetic-site definition,
the complete stalk proof, global sections, the definition of valued
points, and the stated pullback to the ordinary arithmetic spectrum.
This is bounded reading, not a claim to have read the whole paper.

The known programme input is the full, arbitrary bounded distributive
lattice $L$, not a Boolean lattice in its place. Its carrier, complete
prime classification, and supported-zero localization are already proved
in
\href{https://github.com/KokunoYumeto/zeta-function-research-reader/blob/3c78cfbd9e194d35117b33277c421c6564ce3dbf/workbenches/splitzero-tandem/continuations/20260923-full-prime-support/031/01_REBUILT_MATHEMATICS.tex}{031,
Theorem \texttt{thm:rebuilt-full-support}, and Theorem
\texttt{thm:e-stalk}, equation \texttt{eq:e-element-localization}}.
The latter states $G_L(R)[e^{-1}]\simeq L$. It is an existing result,
not a new assertion of the present calculation. We prove its precise
universal consequence for characteristic-one targets, and construct
the resulting full-lattice coefficient extension of the arithmetic
site, including its points, stalks, and Frobenius.

\subsection{The characteristic-one universal property}

Let $R$ be a nonzero commutative unital ring and let $L$ be a nontrivial
bounded distributive lattice, finite or infinite. Retain
\begin{equation}\tag{ASL1}
 \begin{gathered}
 S=G_L(R)=\{(0,\lambda):\lambda\in L\}
                  \cup\{(r,1_L):r\in R\},\\
 (r,\lambda)+(s,\mu)=(r+s,\lambda\vee\mu),\qquad
 (r,\lambda)(s,\mu)=(rs,\lambda\wedge\mu),\\
 z_\lambda=(0,\lambda),\quad \tau=z_{0_L},\quad e=z_{1_L},
 \quad 1_S=(1_R,1_L),\quad
 \sigma:S\longrightarrow L,\quad \sigma(r,\lambda)=\lambda.
 \end{gathered}
\end{equation}
The operations on the semiring $L$ are join and meet, with zero $0_L$
and unit $1_L$. A semiring homomorphism below preserves both zero and
unit.

For every commutative semiring $K$ whose addition is idempotent, the
map $\sigma$ induces a natural bijection
\begin{equation}\tag{ASL2}
 \operatorname{Hom}_{\rm sr}(L,K)
       \xrightarrow{\ \sim\ }
 \operatorname{Hom}_{\rm sr}(S,K),\qquad g\longmapsto g\sigma.
\end{equation}
To prove this, let $f:S\to K$. In $S$, write
$v=(-1_R,1_L)$, so $e=1_S+v$ and $e+1_S=1_S$.
Idempotence in $K$ gives
\[
 f(e)+1_K=(1_K+f(v))+1_K=1_K+f(v)=f(e),
\]
whereas the second source relation makes this equal to $1_K$.
Thus $f(e)=1_K$. Since $e(r,1_L)=e$, every supported amplitude
satisfies $f(r,1_L)=1_K$. Define $g(\lambda)=f(z_\lambda)$.
Addition and multiplication of the $z_\lambda$ show that $g$ preserves
join and meet, and $g(0_L)=0_K$, $g(1_L)=1_K$. It follows that
$f=g\sigma$. Conversely every such composite is a homomorphism.
Surjectivity of $\sigma$ gives uniqueness. This proves (ASL2) directly,
with exactly the same localization map as the established result.

There is a more precise description for an arbitrary idempotent $K$.
Use its natural partial order $x\le y\iff x+y=y$ and let
\begin{equation}\tag{ASL3}
 \operatorname{Id}_{\le1}(K)=\{a\in K:a^2=a,\ a+1_K=1_K\}.
\end{equation}
This is a bounded distributive lattice under addition and multiplication.
For $a,b$ in this set, $ab\le a,b$, while any other $c$ in the set
with $c\le a,b$ satisfies $c=c^2\le ab$. Thus multiplication is
their meet. Also $(a+b)^2=a+b+ab=a+b$, since $ab\le a+b$, so
addition gives the join in the same set. Distributivity is inherited
from $K$. Every unital homomorphism from $L$ has its image here,
because $\lambda^2=\lambda$ and $\lambda\vee1_L=1_L$.
Consequently (ASL2) is also the exact bijection
\begin{equation}\tag{ASL4}
 \operatorname{Hom}_{\rm sr}(S,K)
       \simeq\operatorname{Hom}_{\rm BLat}
                      (L,\operatorname{Id}_{\le1}(K)).
\end{equation}

\subsection{All maps to the actual tropical coefficient semirings}

Connes--Consani define their structure semiring as $\mathbb Z_{\max}$.
Its stalk at the ordered group $H$ is
\[
 H_{\max}=(H\cup\{-\infty\},\max,+).
\]
See their Definition
\texttt{site}, lines 618--622, and Theorem \texttt{structure2},
lines 626--659. Write $\mathbf0=-\infty$ and $\mathbf1=0_H$
in this additive presentation, so that the multiplicative unit is
never confused with the absorbing zero. For any totally ordered
abelian group $H$, the only multiplicative idempotents of $H_{\max}$
are $\mathbf0$ and $\mathbf1$: for $h\in H$, $h+h=h$ forces $h=0_H$.
The same assertion holds for their nonnegative min-version
$(H_{\ge0}\cup\{\infty\},\min,+)$ (source lines 662--665), and
for the multiplicative tropical real semifield
$\mathbb R_+^{\max}=([0,\infty),\max,\cdot)$.

Let $X=\operatorname{SpecLat}L$ be the set of proper lattice prime
ideals. For $P\in X$ set
\begin{equation}\tag{ASL5}
 \chi_P(\lambda)=
 \begin{cases}0,&\lambda\in P,\\1,&\lambda\notin P.
 \end{cases}
 \qquad \chi_P:L\longrightarrow\mathbb B=\{0,1\}.
\end{equation}
A lattice ideal is closed under finite joins and downward closed;
primality means $\lambda\wedge\mu\in P$ forces one factor into $P$.
These properties prove exactly that (ASL5) preserves zero, unit, join,
and meet. Conversely the zero fibre of a bounded lattice homomorphism
to $\mathbb B$ is such a prime ideal, and determines the homomorphism.
Combining this fact with (ASL4) gives, for each of the tropical targets
$K$ just specified,
\begin{equation}\tag{ASL6}
 \operatorname{Hom}_{\rm sr}(S,K)\simeq X,
 \qquad f_P=\iota_{\mathbb B,K}\chi_P\sigma.
\end{equation}
Here the canonical inclusion sends the two Boolean elements to
$\mathbf0_K,\mathbf1_K$. The complete fibres are
\begin{equation}\tag{ASL7}
 \begin{split}
 f_P^{-1}(\mathbf0_K)&=\{z_\lambda:\lambda\in P\},\\
 f_P^{-1}(\mathbf1_K)&=
    \{z_\lambda:\lambda\notin P\}
       \cup\{(r,1_L):r\in R\setminus\{0_R\}\},\\
 f_P^{-1}(a)&=\varnothing\quad
               (a\notin\{\mathbf0_K,\mathbf1_K\}).
 \end{split}
\end{equation}
In particular all ordinary nonzero amplitudes and the supported zero
$e$ have value $\mathbf1_K$. This is forced by the universal property;
it is not an assumption that the added prime has no effect.

The source's Frobenius maps are $x\mapsto x^n$, $n\ge1$, and on
$\mathbb R_+^{\max}$ also $x\mapsto x^t$, $t>0$ (source lines
395--406). They fix both Boolean elements. Thus
\begin{equation}\tag{ASL8}
 \operatorname{Fr}_n^K f_P=f_P,
 \qquad f_P(x^n)=f_P(x),\qquad
 \operatorname{Fr}_t^{\mathbb R_+^{\max}}f_P=f_P.
\end{equation}
The middle identity holds for the actual powers in $S$, by
multiplicativity. It does not assert that $x\mapsto x^n$ is an
additive endomorphism of $S$; that fails, for example, for $R=\mathbb Z$
and $n=2$. On the reflected lattice, $\lambda^n=\lambda$, so its
Frobenius is exactly the identity. Equations (ASL8) are the precise
compatibility statements, with no additional action imposed on $R$.

\subsection{The full lattice is recovered by its whole family of characters}

Define the basic opens $D(\lambda)=\{P\in X:\lambda\notin P\}$.
The map
\begin{equation}\tag{ASL9}
 L\longrightarrow\mathbb B^X,\qquad
 \lambda\longmapsto(\chi_P(\lambda))_{P\in X}
                =\mathbf1_{D(\lambda)}
\end{equation}
is an injective lattice homomorphism. Here and below we use the usual
prime-ideal theorem for bounded distributive lattices, for which the
needed separation proof is as follows. If $\lambda\nleq\mu$, the
ideal $\downarrow\mu$ and filter $\uparrow\lambda$ are disjoint.
By Zorn's lemma choose an ideal maximal among those containing the
first and disjoint from the second. If $a,b$ are outside it but
$a\wedge b$ lies in it, adjoining $a$ and adjoining $b$ each produces
an intersection with the filter. Thus filter elements are bounded
above by $i\vee a$ and $j\vee b$ for some ideal elements $i,j$.
Their meet is bounded above by $i\vee j\vee(a\wedge b)$, by
distributivity, and hence is in the ideal, a contradiction. The
maximal ideal is therefore prime, contains $\mu$, and avoids
$\lambda$. This separates the two elements in (ASL9).

The same proof shows $D(\lambda)$ is compact. If it is covered by
$D(\mu_i)$ but no finite subfamily covers it, then $\lambda$ is not
in the ideal generated by all the $\mu_i$. The preceding ideal/filter
argument gives a prime containing all $\mu_i$ and avoiding $\lambda$,
contradicting the cover. Thus a finite subfamily covers it. Since
\begin{equation}\tag{ASL10}
 D(\lambda\vee\mu)=D(\lambda)\cup D(\mu),\qquad
 D(\lambda\wedge\mu)=D(\lambda)\cap D(\mu),
\end{equation}
every compact open of $X$ equals a unique $D(\lambda)$. Consequently
(ASL9) identifies $L$ with the lattice of compact opens, not with the
entire power set of $X$.

For finite $L$ this is especially explicit. Every nonzero
join-irreducible $j$ is join-prime: if $j\le a\vee b$, then
$j=(j\wedge a)\vee(j\wedge b)$ forces $j\le a$ or $j\le b$.
Its character is $\chi_j(\lambda)=1$ exactly when $j\le\lambda$.
Conversely the complement of any prime ideal is a finite filter;
its least element $j$ is nonzero and join-prime, so every character
arises this way. Thus
\begin{equation}\tag{ASL11}
 X\simeq J(L),\qquad
 L\simeq\{\text{downward-closed subsets of }J(L)\},
 \quad \lambda\longmapsto\{j:j\le\lambda\}.
\end{equation}
For surjectivity in the last map, a downward-closed subset $A$ has
$\lambda=\bigvee A$; join-primality puts every irreducible below
$\lambda$ in $A$. No Boolean support assumption occurs. For the
three-element chain, the character vectors are $(0,0),(1,0),(1,1)$;
the vector $(0,1)$ is absent. Replacing this lattice by two independent
Boolean channels would therefore change the object.

\subsection{Prime centres, localization, and the support structure sheaf}

The complete prime data of the original semiring, already proved in
031, is
\begin{equation}\tag{ASL12}
 \begin{gathered}
 P^{\rm supp}=\{z_\lambda:\lambda\in P\}\quad(P\in X),\\
 Q_{\mathfrak p}=\{z_\lambda:\lambda\in L\}
              \cup\{(r,1_L):r\in\mathfrak p\}
                    \quad(\mathfrak p\in\operatorname{Spec}R),\\
 \operatorname{Spec}S=X\blacktriangleright\operatorname{Spec}R,
 \qquad D_S(e)=X,\quad V_S(e)=\operatorname{Spec}R.
 \end{gathered}
\end{equation}
For clarity, exhaustion follows directly from the source operations.
An ideal containing a supported element contains $e$, hence every
$z_\lambda$; its supported amplitudes form a ring ideal. An ideal
containing no supported element consists of the $z_\lambda$ from a
proper lattice ideal. The two primality tests are respectively ring
primality and meet-primality. Every support prime lies strictly below
every arithmetic prime, and the basic opens are
$D_S(z_\lambda)=D(\lambda)$ and
$D_S((r,1_L))=X\cup D_R(r)$. This proves the stated topology.

A tropical semifield has only its zero ideal as a proper ideal,
because every nonzero element is a unit. The spectral map induced by
$f_P$ therefore has the single image $P^{\rm supp}$, by (ASL7).
More generally, every proper prime of any target in (ASL6) contracts
to that same prime, since it contains $\mathbf0$ and excludes
$\mathbf1$. In particular these maps do not centre on an arithmetic
prime $Q_{\mathfrak p}$ containing $e$.

For a multiplicative subset $T\subset S$, a character map $f_P$
extends to $T^{-1}S$ exactly when
\begin{equation}\tag{ASL13}
 \sigma(T)\cap P=\varnothing;
 \qquad f_P(s/t)=f_P(s)\quad(t\in T).
\end{equation}
Indeed all target values of denominators are either zero or unit one.
For the universal property with arbitrary idempotent targets,
localization commutes with the reflection:
\begin{equation}\tag{ASL14}
 (T^{-1}S)_{\rm idemp}\simeq\sigma(T)^{-1}L.
\end{equation}
Both sides represent maps from $S$ to idempotent semirings in which
all elements of $T$ become units, using (ASL2), so the displayed
isomorphism follows from the two inverse universal maps.

The full support localizations are explicit:
\begin{equation}\tag{ASL15}
 L[\lambda^{-1}]\simeq\downarrow\lambda,
 \qquad \mu\longmapsto\mu\wedge\lambda,
 \quad 1_{\downarrow\lambda}=\lambda.
\end{equation}
The equivalence relation is
$\mu\sim\nu\iff\mu\wedge\lambda=\nu\wedge\lambda$;
all positive powers of the denominator are $\lambda$ itself.
For a prime $P$, the stalk has the description
\begin{equation}\tag{ASL16}
 L_P=L/(\mu\sim_P\nu),\qquad
 \mu\sim_P\nu\iff
   \exists t\notin P:\mu\wedge t=\nu\wedge t,
 \qquad S_{P^{\rm supp}}\simeq L_P.
\end{equation}
Every denominator in $L\setminus P$ is idempotent, so after being
inverted it becomes the unit; this gives the stated equivalence and
exhausts fractions. In $S_{P^{\rm supp}}$, $e$ is inverted first,
giving $L$, and the remaining denominators are precisely those
represented by $L\setminus P$. This proves the last isomorphism.

These semirings of sections form a sheaf on the
support spectrum. On compact basic opens set
\begin{equation}\tag{ASL17}
 \mathcal L(D(\lambda))=\downarrow\lambda,
 \qquad
 \operatorname{res}_{D(\mu)}(a)=a\wedge\mu
                   \quad(\mu\le\lambda).
\end{equation}
For a finite cover with $\lambda=\bigvee_i\lambda_i$, compatible
sections $a_i\le\lambda_i$ obey
$a_i\wedge\lambda_j=a_j\wedge\lambda_i$. Their unique gluing is
$a=\bigvee_i a_i$: distributivity proves
$a\wedge\lambda_i=a_i$, and if a second section has these restrictions
then intersecting it with $\bigvee_i\lambda_i$ proves equality.
Compactness gives a finite subcover for any cover of a basic open,
and the same compatibility checks every remaining member. Thus this
is a sheaf on the basis and has stalk (ASL16). It is the restriction
of the localization structure sheaf of $S$ to $D_S(e)$.

The stalk $L_P$ generally retains more than a Boolean value. For the
chain $0<a<1$ at the prime $P=\{0,a\}$, only $1$ is inverted, so
$L_P=L$ still has three elements. Its local residue map to
$\mathbb B$ sends $0,a$ to zero and $1$ to one. In general the
extension of $\chi_P$ sends a class to one precisely when that class
is the unit: $\mu\notin P$ becomes invertible and hence equals one;
if $\mu\in P$ could equal one, then some $t\notin P$ would satisfy
$t\le\mu$, contradicting downward closure. The Bourne quotient
by its zero fibre is exactly $\mathbb B$, because it collapses all
nonunits to zero while leaving the unit distinct. Thus the residue
character is related to the full support stalk by an exact local
map; it is not a replacement for that stalk.

\subsection{The universal tropical coefficient extension retains every support}

Let $K$ be a tropical semifield as above and form the commutative
semiring tensor product
\begin{equation}\tag{ASL18}
 K_L=K\otimes_{\mathbb B}L.
\end{equation}
There is no unital map $\mathbb B\to S$: it would require
$1_S+1_S=1_S$, which in the nonzero ring amplitude would say
$2\cdot1_R=1_R$. The correct direct base change from $S$ uses the
initial semiring $\mathbb N_0$, and (ASL2) gives the canonical
isomorphism
\begin{equation}\tag{ASL19}
 S\otimes_{\mathbb N_0}K\simeq L\otimes_{\mathbb B}K.
\end{equation}
Indeed any semiring receiving a unital map from $K$ is additively
idempotent: $1+1=1$ implies $x+x=x$ for every element. Its map from
$S$ therefore factors uniquely through $L$. Maps from $L$ and $K$
agree on $\mathbb B$, and these are exactly the universal maps on
the right side of (ASL19). This proves the isomorphism without
pretending that $S$ was already a Boolean algebra or a
$\mathbb B$-algebra.

The simultaneous character evaluations give a faithful model
\begin{equation}\tag{ASL20}
 \begin{split}
 K_L&\longrightarrow K^X,\\
 w=\sum_{i=1}^r a_i\otimes\lambda_i
    &\longmapsto\left(P\longmapsto
          \max_{i:\lambda_i\notin P}a_i\right),
 \end{split}
\end{equation}
with the empty maximum equal to $\mathbf0_K$.
Both operations are pointwise on the displayed functions. The map
exists by the tensor universal property. Every tensor expression is
a finite sum of pure tensors, because products of pure tensors are
pure tensors. To prove injectivity, discard zero-coefficient terms
and list the finitely many distinct remaining coefficients as
$a_1<\cdots<a_r$ in the natural total order of $K$. Let
$\mu_i$ be the join of the supports whose coefficients are at least
$a_i$. The original tensor expression equals
\begin{equation}\tag{ASL21}
 \sum_{i=1}^r a_i\otimes\mu_i,
 \qquad \mu_1\ge\cdots\ge\mu_r.
\end{equation}
For verification, distributing each $\mu_i$ merely adds terms
$a_i\otimes\lambda$ already bounded above by a term
$a_j\otimes\lambda$ with $j\ge i$; their sum is the latter because
$a_i+a_j=a_j$. Thus all original data is related by actual tensor
relations. In the image function $w(P)$, the set where $w(P)\ge a_i$
is exactly $D(\mu_i)$. For two expressions, use the union of their
finite coefficient lists in (ASL21). Equality of functions gives
equality of every $D(\mu_i)$ and therefore of every $\mu_i$ by
(ASL9). Hence the two tensors are equal, proving faithfulness.

Equivalently, (ASL20) identifies $K_L$ with the finite-valued
functions $w:X\to K$ whose sets $\{P:w(P)\ge a\}$ are compact open
for every nonzero attained value $a$. Conversely, the finitely many
nonzero values of such a function and their compact opens supply
(ASL21), by (ASL10). In particular $1\otimes L$ embeds faithfully
as its compact-open characteristic functions. The fibre of evaluation
at a specified $P$ is completely described by the displayed maximum:
it is zero exactly when all nonzero-coefficient supports belong to
$P$; a nonzero value $a$ occurs exactly when $a$ is the largest active
coefficient. In level-set language this is
$P\in D(\mu_{\ge a})\setminus D(\mu_{>a})$.

The coefficient extension can also accompany the original ring
amplitudes without losing them. The map
\begin{equation}\tag{ASL22}
 S\longrightarrow R\times K_L,\qquad
 (r,\lambda)\longmapsto(r,1\otimes\lambda)
\end{equation}
is an injective semiring homomorphism: its first coordinate retains
the amplitude and (ASL20) makes its second coordinate faithful on
the entire lattice. Its image consists exactly of pairs with a
compact-open characteristic function in the second coordinate and
with that function identically one whenever $r\ne0$. This is an
actual correspondence retaining both coordinates, not a claim that
the ring amplitude itself embeds into a characteristic-one semiring.

\subsection{The arithmetic site with full lattice coefficients}

Let $\mathscr A=\widehat{\mathbb N^\times}$ be the source topos and
$\mathcal O=\mathbb Z_{\max}$ with its Frobenius action, exactly as
in Connes--Consani's Definition \texttt{site}. Give the constant
semiring object $\underline L$ the trivial action and set
\begin{equation}\tag{ASL23}
 \mathcal O_L=\mathcal O\otimes_{\underline{\mathbb B}}\underline L,
 \qquad
 (\mathscr A,\mathcal O_L)\longrightarrow(\mathscr A,\mathcal O).
\end{equation}
The underlying geometric morphism is the identity, and its structure
map is $\mathcal O\to\mathcal O_L$, $a\mapsto a\otimes1$.
By (ASL19) it is also the universal coefficient object receiving
$\mathcal O$ and the constant full semiring $\underline S$ over
$\underline{\mathbb N_0}$. The latter is expressly a constant object;
no unproved additive Frobenius on the ring $R$ is asserted.

The Frobenius action and joint character maps are
\begin{equation}\tag{ASL24}
 \operatorname{Fr}_n\left(\sum_i a_i\otimes\lambda_i\right)
       =\sum_i a_i^n\otimes\lambda_i,
 \qquad
 \operatorname{ev}_P:\mathcal O_L\to\mathcal O.
\end{equation}
Evaluation is (ASL20), so these are equivariant semiring maps and
their family is jointly injective. In the additive notation of
$\mathbb Z_{\max}$ the coefficient operation $a_i^n$ means $na_i$;
the absorbing zero stays absorbing zero. For each $P$ the exact
coefficient specialization is
\begin{equation}\tag{ASL25}
 \mathcal O_L\otimes_{\underline L,\chi_P}
           \underline{\mathbb B}\simeq\mathcal O.
\end{equation}
This follows from the two tensor universal properties, or directly
by substituting $\lambda\mapsto\chi_P(\lambda)$ in every pure
tensor. All supports are present before this specified specialization.

At the point corresponding to a nonzero ordered subgroup
$H\subset\mathbb Q$, the stalk is
\begin{equation}\tag{ASL26}
 (\mathcal O_L)_H\simeq H_{\max}\otimes_{\mathbb B}L.
\end{equation}
To justify this precisely, the inverse-image functor of the point
preserves finite products and arbitrary colimits. A semiring tensor
product is constructed from finite products, coproducts and the
coequalizers imposing its addition and multiplication relations.
It is therefore preserved by that inverse-image functor. A constant
object with trivial action has the same stalk: in the source's
description $(X\times H_{>0})/\sim$ at lines 631--649, any two
positive parameters have a common divisor in $H_{>0}$, and the
trivial action identifies them while retaining the value in $X$.
Finally the source's Theorem \texttt{structure2} identifies the
original structure stalk as $H_{\max}$. These facts prove (ASL26).
Its elements have exactly the full-lattice function description
(ASL20)--(ASL21), now with coefficients in $H_{\max}$.

The global sections of the extended site are exactly
\begin{equation}\tag{ASL27}
 \Gamma(\mathscr A,\mathcal O_L)\simeq L.
\end{equation}
For a one-object action topos, a global section is a simultaneous
fixed element; this is the source's argument at lines 668--673.
Under the faithful function representation (ASL20), invariance under
$\operatorname{Fr}_2$ forces each value $a$ to satisfy $a^2=a$.
In $\mathbb Z_{\max}$ these are precisely $\mathbf0,\mathbf1$.
Such a finite-valued function is a section of $\mathcal O_L$ exactly
when its nonzero set is $D(\lambda)$, and then it is
$1\otimes\lambda$. Conversely these sections are fixed by every
Frobenius. This proves (ASL27); the source's unextended case is
$\Gamma(\mathscr A,\mathcal O)=\mathbb B$.

\subsection{All tropical-valued points of the extended site}

By the definition of a valued point in the source, lines 685--690,
one supplies a topos point $H$ and a semiring map from its structure
stalk to $\mathbb R_+^{\max}$. The tensor universal property and
(ASL6) therefore give the exact classification
\begin{equation}\tag{ASL28}
 \begin{split}
 \operatorname{Hom}_{\rm sr}
   (H_{\max}\otimes_{\mathbb B}L,\mathbb R_+^{\max})
 &\simeq
 \operatorname{Hom}_{\rm sr}(H_{\max},\mathbb R_+^{\max})\times X,\\
 (\beta,P):\quad
 \sum_i a_i\otimes\lambda_i
 &\longmapsto\max_{i:\lambda_i\notin P}\beta(a_i).
 \end{split}
\end{equation}
There is no missing compatibility condition: every unital semiring
map agrees on the two distinguished elements of $\mathbb B$.
Thus the set of isomorphism classes of valued points of the extended
site is the product of the original valued-point set with $X$.
An isomorphism of underlying topos points can change $H$ and its
map $\beta$ as in the source's equivalence relation, but it acts
identically on the constant object $L$, hence cannot identify
different $P$'s.

For a fixed embedding $H\subset\mathbb Q$, the first factor in
(ASL28) can itself be given exactly. Every map has the form
\begin{equation}\tag{ASL29}
 \beta_c(-\infty)=0,\qquad
 \beta_c(h)=e^{ch}\ (h\in H),\qquad c\ge0.
\end{equation}
Indeed group units must map to positive real units. Taking their
real logarithms gives an additive order-preserving map
$\ell:H\to\mathbb R$. Fix $h_0>0$. Since every $h/h_0$ is rational,
additivity forces $\ell(h)=h\ell(h_0)/h_0$; order preservation makes
$c=\ell(h_0)/h_0\ge0$. Conversely every formula (ASL29) preserves
both multiplication and maximum. The value $c=0$ is the degenerate
map through $\mathbb B$; $c>0$ is injective on the ordered group.

Under the source's real Frobenius the exact point action is
\begin{equation}\tag{ASL30}
 \operatorname{Fr}_t:(H,\beta_c,P)\longmapsto
                         (H,\beta_{tc},P),\qquad t>0.
\end{equation}
Thus the arithmetic site's original nontrivial coefficient dynamics
is retained, and the full lattice contributes its whole space of
prime characters as a fixed parameter. A direct map from $S$ to a
single tropical stalk sees only its one character, by (ASL7), whereas
the extended coefficient sheaf, all its stalks, and their joint
evaluations retain every element and every relation of $L$.
Equations (ASL19), (ASL23), (ASL26), and (ASL28) are the exact
connecting maps. No purity, cohomology comparison, or Riemann-hypothesis
conclusion is inferred from them.

\subsection{A precise distinction between prime ideals and valued points}

The preceding statements classify all support primes of $S$, all
lattice primes, and all tropical-valued points of the coefficient
extension. They do not identify the ordinary semiring prime spectrum
of $K_L$ with $X$. There can be additional primes, with an explicit
example already for the three-element chain $L=\{0<a<1\}$.
For $K=\mathbb R_+^{\max}$, (ASL20) gives
\begin{equation}\tag{ASL31}
 K_L\simeq\{(u,v)\in[0,\infty)^2:u\ge v\},\qquad
 \mathfrak q=\{(0,0)\}\cup\{(u,v):u>v\}.
\end{equation}
The operations are coordinatewise maximum and multiplication.
The set $\mathfrak q$ is a proper prime ideal. It is closed under
maximum: whichever of two first coordinates is larger also exceeds
both second coordinates. Multiplying $(u,v)$ with $u>v$ by a
nonzero $(r,s)$ with $r\ge s$ gives $ur>vs$; multiplication by zero
gives zero. Its complement is the set of positive diagonal pairs,
which is multiplicatively closed. These observations prove the ideal
and prime assertions. But $\mathfrak q$ is neither evaluation kernel:
the first evaluation has only the zero kernel, the second kernel is
$\{(u,0):u\ge0\}$, and $(2,1)$ belongs to $\mathfrak q$ but to
neither of those. Thus the full support and valued-point comparison
retains its precise stated scope even in this smallest non-Boolean
case. It does not erase other prime geometry introduced by the
tropical coefficient extension.

\subsection{The exact pullback over the ordinary arithmetic spectrum}

Connes--Consani's geometric morphism over $\operatorname{Spec}\mathbb Z$
uses the monoid $\mathbb N_0^\times$ that includes zero, not just
$\mathbb N^\times$. This distinction is explicit in their Remark
\texttt{caseno}, source line 660, and Theorem \texttt{thmspz}, lines
951--961. Put $\mathscr A_0=\widehat{\mathbb N_0^\times}$ and let
$\mathcal O^0$ be their extended structure semiring. In additive
exponent notation its additional action is
\begin{equation}\tag{ASL32}
 \operatorname{Fr}_0(-\infty)=-\infty,\qquad
 \operatorname{Fr}_0(h)=0_H\quad(h\in H).
\end{equation}
This preserves both operations: a maximum is absorbing zero exactly
when both inputs are absorbing zero, and a product is absorbing zero
exactly when one factor is. The extension
$\mathcal O^0_L=\mathcal O^0\otimes_{\underline{\mathbb B}}
\underline L$ has
\begin{equation}\tag{ASL33}
 \operatorname{Fr}_0\!\left(\sum_i a_i\otimes\lambda_i\right)
 =1\otimes\rho\!\left(\sum_i a_i\otimes\lambda_i\right),
 \qquad
 \rho\!\left(\sum_i a_i\otimes\lambda_i\right)
       =\bigvee_{a_i\ne\mathbf0}\lambda_i.
\end{equation}
The map $\rho$ is well defined by the tensor universal property,
using the semiring map $H_{\max}\to\mathbb B$ that sends every
nonzero element to one. Equivalently, in (ASL20), $D(\rho(w))$
is exactly the set where $w$ is nonzero. Thus (ASL33) retains a
specific element of $L$, rather than only whether any support exists.
It obeys $\operatorname{Fr}_0\operatorname{Fr}_n
=\operatorname{Fr}_n\operatorname{Fr}_0=\operatorname{Fr}_0$.

The original source constructs
$\Theta:\operatorname{Sh}(\operatorname{Spec}\mathbb Z)
\to\mathscr A_0$. At a rational prime $p$, its point is
$H(p)=\mathbb Z[1/p]\subset\mathbb Q$, as defined at source lines
910--914; the generic point $\eta$ goes to the trivial group.
Applying its inverse-image functor to the full coefficient extension
gives the exact sheaf and stalk formulas
\begin{equation}\tag{ASL34}
 \begin{gathered}
 \mathcal T_L:=\Theta^*\mathcal O^0_L
   \simeq\Theta^*\mathcal O^0\otimes_{\underline{\mathbb B}}
                \underline L,\\
 (\mathcal T_L)_p\simeq
      \mathbb Z[1/p]_{\max}\otimes_{\mathbb B}L,
 \qquad (\mathcal T_L)_\eta\simeq L.
 \end{gathered}
\end{equation}
The preservation argument proving (ASL26) applies to $\Theta^*$ as
well. Connes--Consani's Proposition \texttt{sheavesonspz0}, source
lines 965--975, gives the original stalks $H(p)_{\max}$ and
$\mathbb B$, so tensoring proves (ASL34). This is a pullback of
semiring coefficients on the underlying topos of the arithmetic
spectrum. It does not identify these coefficients with the ordinary
ring structure sheaf $\mathcal O_{\operatorname{Spec}\mathbb Z}$.

All its sections can also be described without suppressing the prime
or support labels. For a nonempty open $U$ of
$\operatorname{Spec}\mathbb Z$, a section of $\mathcal T_L$ is
exactly the data
\begin{equation}\tag{ASL35}
 \begin{gathered}
 \lambda\in L,\qquad
 w_p\in\mathbb Z[1/p]_{\max}\otimes_{\mathbb B}L
                  \quad(p\in U\text{ a rational prime}),\\
 \rho(w_p)=\lambda\quad\hbox{for every }p,\qquad
 w_p=1\otimes\lambda\quad\hbox{for all but finitely many }p.
 \end{gathered}
\end{equation}
The value at $\eta$ is $\lambda$. Addition and multiplication are
pointwise, with join and meet at $\eta$; restriction discards the
coordinates at the removed primes. The integer Frobenius applies
$h\mapsto nh$ to every coefficient at every retained prime and
acts identically on $\lambda$. Its zero action is (ASL33).

Here is a complete gluing proof. The source's Proposition
\texttt{sheavesonspz0}(ii) says that an original section is either
absorbing zero everywhere or a family of nonzero coefficients
equal to the unit outside a finite set of primes. A local finite
tensor expression therefore has a single generic support
$\lambda=\bigvee_{a_i\ne\mathbf0}\lambda_i$; all its prime values
have that same support, and outside a finite set every value is
$1\otimes\lambda$. Every sheaf tensor section has such expressions
locally, by the construction of the tensor as the sheafification of
the presheaf tensor. Every nonempty open $U$ here is quasi-compact:
a nonempty member of any open cover misses only finitely many points
of $U$, which finitely many other members can cover. Thus finitely
many local expressions suffice. Their generic values agree since
every nonempty intersection contains $\eta$. Their finite exceptional
sets have finite union. This proves that every section yields (ASL35).

Conversely take (ASL35), with finite exceptional set $E$. For each
$p\in E$ write $w_p$ as a finite tensor sum using nonzero
coefficients $a_i\in\mathbb Z[1/p]_{\max}$. Its support labels
have join $\lambda$. On the open $U_p=U\setminus(E\setminus\{p\})$
extend $a_i$ to the original coefficient section with value $a_i$
at $p$ and the unit at every other prime; the source's section
description makes this an actual section. Tensor these sections
with their labels and add. The result has prescribed value $w_p$
at $p$ and value $1\otimes\lambda$ elsewhere on $U_p$.
The case $\lambda=0_L$ has only zero tensors and is obtained from
the zero section. Add the constant section $1\otimes\lambda$ on
$U\setminus E$. These local sections agree on their intersections,
so the sheaf gluing axiom constructs the required global section.
Uniqueness follows because sections of a sheaf on a topological
space are determined by all their germs. This proves (ASL35) and
its operation and restriction formulas.
